\documentclass[a4paper,12pt]{article}
\usepackage[T2A]{fontenc}
\usepackage[utf8]{inputenc}
\usepackage[russian]{babel}
\usepackage{amsmath,amssymb,mathtools}
\usepackage{helvet}
\renewcommand{\familydefault}{\sfdefault}
\usepackage{icomma}
\usepackage[a4paper,left=22mm,right=22mm,top=20mm,bottom=20mm,headheight=14pt,headsep=8mm]{geometry}
\usepackage{microtype}
\usepackage{tikz}
\usetikzlibrary{arrows.meta,positioning,shapes.geometric,calc}
\usepackage{booktabs,tabularx}
\usepackage{enumitem}
\usepackage[hidelinks]{hyperref}
\usepackage{parskip}
\usepackage{fancyhdr}
\usepackage{setspace}
\usepackage{xcolor}
\definecolor{accent}{HTML}{008080}
\definecolor{darkaccent}{HTML}{004D4D}
\definecolor{textgray}{HTML}{333333}
\definecolor{mute}{HTML}{67747A}
\color{textgray}
\setstretch{1.36}
\setlength{\parindent}{0pt}
\setlength{\parskip}{5pt}
\setlist[itemize]{leftmargin=*,topsep=2pt,itemsep=1pt}
\setlist[enumerate]{leftmargin=*,topsep=4pt,itemsep=6pt}
\pagestyle{fancy}
\fancyhf{}
\fancyhead[L]{\small\color{mute}Движение вдогонку и задачи на движение}
\fancyhead[R]{\small\color{mute}Ярослав Верещагин}
\fancyfoot[C]{\small\color{mute}\thepage}
\renewcommand{\headrulewidth}{0pt}
\renewcommand{\footrulewidth}{0pt}
\newcommand{\topic}[1]{\section*{\color{darkaccent}#1}}
\newcommand{\subtopic}[1]{\subsection*{\color{accent}#1}}
\newcommand{\key}[1]{\textbf{\color{darkaccent}#1}}
\newcommand{\kmh}{\text{ км/ч}}
\newcommand{\kms}{\text{ км}}
\newcommand{\minu}{\text{ мин}}
\newcommand{\hour}{\text{ ч}}
\newcommand{\met}{\text{ м}}
\newcommand{\tabheaders}{\toprule Участник движения & Скорость, км/ч & Время, ч & Путь, км\\\midrule}
\begin{document}
\thispagestyle{empty}
\begin{center}
\vspace*{20mm}
{\large\color{darkaccent}МАТЕМАТИКА}\par
\vspace{18mm}
{\LARGE Чек-лист}\par
\vspace{12mm}
{\LARGE\color{accent}Движение вдогонку\par и задачи на движение}\par
\vspace{9mm}
{\large Схемы, таблицы, уравнения и единицы измерения}\par
\vspace{10mm}
{\normalsize Подготовка к ОГЭ по математике}\par
\vfill
{\large Ярослав Верещагин}\par
\vspace{10mm}
5 ноября 2026 г.\par
\vspace{15mm}
{\small Лёвин Артём Александрович\par эксклюзивно для Ярослава Верещагина}\par
\vspace{8mm}
\begin{tikzpicture}
\draw[accent,line width=0.8pt] (0,0) .. controls (1.1,0.7) and (2.1,-0.5) .. (3.2,0)
 .. controls (4.4,0.5) and (5.3,-0.6) .. (6.4,0);
\fill[darkaccent] (3.2,0) circle (1pt);
\end{tikzpicture}
\vspace*{8mm}
\end{center}
\clearpage

\topic{1. Базовая модель движения}
В текстовых задачах важно сначала выяснить, \key{кто движется, откуда, куда и в какой момент}. После этого вводят обозначения и составляют таблицу. Самое важное равенство:
\[
\boxed{S=vt},\qquad t=\frac{S}{v},\qquad v=\frac{S}{t}.
\]
Здесь $S$ обозначает пройденный путь, $v$ --- скорость движения, а $t$ --- время в пути.

\key{Единицы должны быть согласованы.} Если скорость записана в километрах в час, расстояние удобно выражать в километрах, а время --- в часах. В итоговом ответе указывают именно те единицы, которые требует задача.
\[
200\met=0{,}2\kms,\qquad 8\minu=\frac{8}{60}\hour=\frac{2}{15}\hour.
\]

\subtopic{Три вопроса перед составлением уравнения}
\begin{enumerate}
\item \textbf{Откуда началось движение?} Участники вышли из одного места или находились на расстоянии друг от друга?
\item \textbf{Как связаны пути?} Один прошёл дальше, оба прошли одинаковое расстояние или один успел развернуться?
\item \textbf{Как связаны времена?} Двигались одинаковое время, кто-то вышел раньше или одному потребовалось больше времени?
\end{enumerate}

\key{Полезная привычка:} после каждой полученной строки проверяйте \textbf{скобки, знаки и размерности}. Именно эти три проверки помогают заметить большинство вычислительных ошибок.

\topic{2. Два участника одновременно идут в одном направлении}
Если люди начали движение из одного места и одновременно, то к одному и тому же моменту их времена одинаковы. Быстрый участник проходит большее расстояние. Разность пройденных путей равна расстоянию между ними.
\[
S_{\text{быстр}}-S_{\text{медл}}=d,
\qquad (v_{\text{быстр}}-v_{\text{медл}})\,t=d.
\]
Здесь $d$ --- расстояние между участниками движения в рассматриваемый момент.

\subtopic{Разобранный пример: расстояние между пешеходами}
Два пешехода одновременно отправились из одного места в одном направлении. Скорость первого на $1\kmh$ больше скорости второго. Через сколько минут расстояние между пешеходами станет равным $200\met$?

Обозначим скорость медленного пешехода через $y\kmh$, а время движения через $x\hour$. Тогда скорость быстрого равна $(y+1)\kmh$.
\begin{center}
\renewcommand{\arraystretch}{1.2}
\begin{tabularx}{\linewidth}{Xccc}
\tabheaders
Быстрый пешеход & $y+1$ & $x$ & $x(y+1)$\\
Медленный пешеход & $y$ & $x$ & $xy$\\
\bottomrule
\end{tabularx}
\end{center}
Переведём расстояние в километры: $200\met=0{,}2\kms$. Вычтем из большего пути меньший:
\[
x(y+1)-xy=0{,}2,\qquad xy+x-xy=0{,}2,\qquad x=0{,}2\hour.
\]
Так как $1\hour=60\minu$, получаем $0{,}2\cdot60=12\minu$.
\par\key{Ответ:} через $12$ минут.

\topic{3. Один участник разворачивается и идёт обратно}
Если быстрый человек дошёл до конечного пункта, развернулся и встретил медленного на обратном пути, \key{расстояния, пройденные ими от старта, различны}. Для записи путей нужен рисунок.

\begin{center}
\begin{tikzpicture}[x=1cm,y=1cm,>=Stealth]
\coordinate (A) at (0,0); \coordinate (M) at (5.9,0); \coordinate (B) at (9.3,0);
\draw[line width=1.0pt] (A)--(9.9,0);
\foreach \p in {A,M,B}{\fill[darkaccent] (\p) circle (2pt);}
\node[below=5pt] at (A) {Старт $A$};
\node[below=5pt] at (M) {Встреча $M$};
\node[below=5pt] at (B) {Опушка $B$};
\draw[->,line width=1.1pt,accent] (0,0.46)--(9.18,0.46);
\draw[->,line width=1.1pt,accent] (9.3,0.65)--(5.98,0.65);
\draw[->,line width=1.1pt,orange!80!black] (0,-0.7)--(5.85,-0.7);
\node[above=2pt,accent] at (4.3,0.49) {Быстрый: до опушки и обратно};
\node[below=2pt,orange!80!black] at (2.9,-0.7) {Медленный: до места встречи};
\draw[<->,line width=0.8pt] (5.9,-1.36)--node[fill=white,inner sep=1pt]{$x$}(9.3,-1.36);
\end{tikzpicture}
\end{center}

Если от старта до опушки $L$, а от опушки до места встречи $x$, быстрый проходит $L+x$, медленный --- $L-x$. Начали они одновременно, поэтому времена их движения равны:
\[
\frac{L+x}{v_{\text{быстр}}}=\frac{L-x}{v_{\text{медл}}}.
\]
\textbf{Важно:} $0<x<L$, так как встреча происходит между стартом и опушкой, после того как быстрый пешеход успел развернуться.

\subtopic{Разобранный пример: прогулка до опушки леса}
Два человека одновременно отправились из одного места к опушке леса, до которой $6\kms$. Первый идёт со скоростью $6{,}5\kmh$, а второй --- со скоростью $3{,}5\kmh$. Первый дошёл до опушки и повернул обратно. На каком расстоянии от опушки они встретились?

Обозначим искомое расстояние через $x\kms$. Быстрый прошёл $6+x$ километров, медленный --- $6-x$ километров. Их времена равны:
\[
\frac{6+x}{6{,}5}=\frac{6-x}{3{,}5}.
\]
Перемножим числитель первой дроби со знаменателем второй и наоборот:
\[
3{,}5(6+x)=6{,}5(6-x).
\]
Умножим обе части на $2$, чтобы избавиться от десятичных дробей:
\[
7(6+x)=13(6-x),\qquad 42+7x=78-13x.
\]
\[
20x=36,\qquad x=1{,}8\kms=1800\met.
\]
\key{Ответ:} $1800$ метров от опушки.

\topic{4. Один участник выехал раньше другого}
При движении по одному и тому же маршруту длина пути может быть одинаковой, хотя скорости различаются. В таком случае удобно выразить время каждого участника:
\[
t_{1}=\frac{S}{v_{1}},\qquad t_{2}=\frac{S}{v_{2}}.
\]
Если более медленный участник выехал раньше, а оба прибыли одновременно, \key{его время движения больше}. Разность времён равна промежутку между отправлениями:
\[
\frac{S}{v_{\text{медл}}}-\frac{S}{v_{\text{быстр}}}=\Delta t.
\]

\subtopic{Разобранный пример: грузовик и автобус}
Из посёлка $A$ в посёлок $B$, расстояние между которыми $20\kms$, выехал грузовик. Через $8\minu$ следом за ним отправился автобус, скорость которого на $5\kmh$ больше скорости грузовика. Оба автомобиля прибыли в посёлок $B$ одновременно. Найдите скорость автобуса.

Пусть скорость автобуса равна $x\kmh$. Тогда скорость грузовика равна $(x-5)\kmh$. Так как обе скорости должны быть положительными, $x>5$.
\begin{center}
\renewcommand{\arraystretch}{1.2}
\begin{tabularx}{\linewidth}{Xccc}
\tabheaders
Грузовик & $x-5$ & $\dfrac{20}{x-5}$ & $20$\\[5pt]
Автобус & $x$ & $\dfrac{20}{x}$ & $20$\\[5pt]
\bottomrule
\end{tabularx}
\end{center}
Грузовик двигался на $8\minu$ дольше. Переведём минуты в часы:
\[
\frac{20}{x-5}-\frac{20}{x}=\frac{8}{60}=\frac{2}{15}.
\]
Сократим коэффициенты в уравнении, разделив обе части на $2$:
\[
\frac{10}{x-5}-\frac{10}{x}=\frac{1}{15}.
\]
Приведём дроби слева к общему знаменателю:
\[
\frac{10x-10(x-5)}{x(x-5)}=\frac{1}{15},\qquad
\frac{50}{x(x-5)}=\frac{1}{15}.
\]
При $x>5$ можно перемножить крест-накрест:
\[
x(x-5)=750,\qquad x^2-5x-750=0.
\]
Разложим на множители: $(x-30)(x+25)=0$. Корни $x=30$ и $x=-25$. Отрицательная скорость невозможна, поэтому $x=30$.

\key{Проверка:} грузовик двигался $20/25=0{,}8$ часа, автобус --- $20/30=\frac23$ часа. Разность равна $\frac{2}{15}$ часа, то есть $8$ минут.

\key{Ответ:} $30\kmh$.

\topic{5. Разность скоростей в непривычных единицах}
Иногда в условии сравнивают скорости в метрах в минуту, а путь задают в километрах. Прежде чем заполнять таблицу, нужно перевести \key{скорость как отношение двух величин}.
\[
200\frac{\text{м}}{\text{мин}}
=200\cdot\frac{1\text{ км}}{1000\text{ м}}\cdot\frac{60\text{ мин}}{1\text{ ч}}
=12\kmh.
\]
Мы делим на $1000$, чтобы перейти от метров к километрам, и умножаем на $60$, чтобы перейти от величины «за минуту» к величине «за час».

\subtopic{Разобранный пример: пассажирский и скорый поезда}
Пассажирский поезд за каждую минуту проходит на $200\met$ меньше, чем скорый. На путь в $360\kms$ пассажирский поезд затратил на один час больше скорого. Найдите скорость пассажирского поезда.

Пусть скорость пассажирского поезда $x\kmh$. Скорость скорого тогда $x+12\kmh$. Поскольку пассажирский поезд ехал дольше, имеем:
\[
\frac{360}{x}-\frac{360}{x+12}=1,\qquad x>0.
\]
Приведём дроби к общему знаменателю:
\[
\frac{360(x+12)-360x}{x(x+12)}=1,\qquad
x^2+12x=4320.
\]
\[
x^2+12x-4320=0,\qquad (x-60)(x+72)=0.
\]
Положительный корень $x=60$. Скорость скорого поезда равна $72\kmh$.

\key{Проверка:} $360/60=6$ часов, $360/72=5$ часов. Разница действительно равна одному часу.

\key{Ответ:} $60\kmh$.

\subtopic{Ещё один перевод единиц}
\[
1\kms=1000\met=100000\text{ см},\qquad
15\frac{\text{см}}{\text{ч}}=\frac{15}{100000}\kmh=0{,}00015\kmh.
\]
Число также можно записать в стандартном виде: $1{,}5\cdot10^{-4}\kmh$. Обратите внимание, что преобразуется только единица длины: часы уже совпадают.

\clearpage
\topic{6. Алгоритм решения текстовой задачи на движение}
\begin{center}
\begin{tikzpicture}[>=Stealth, node distance=9mm,
 startstop/.style={ellipse,draw=darkaccent,line width=0.9pt,minimum width=47mm,minimum height=8mm,align=center,font=\small},
 process/.style={rounded corners=2pt,draw=accent,line width=0.8pt,text width=115mm,minimum height=13mm,align=center,font=\small,inner sep=5pt},
 arrow/.style={->,line width=0.9pt,draw=darkaccent}]
\node[startstop] (a) {Начало};
\node[process,below=of a] (b) {Прочитайте условие. Определите участников, маршрут, момент старта и вопрос задачи.};
\node[process,below=of b] (c) {Нарисуйте схему движения, если по ней легче увидеть пути и место встречи.};
\node[process,below=of c] (d) {Обозначьте неизвестные величины буквами и согласуйте единицы измерения.};
\node[process,below=of d] (e) {Заполните таблицу: скорость, время, расстояние. Используйте $S=vt$.};
\node[process,below=of e] (f) {Выберите связь по условию: равенство времён, разность времён или разность пройденных путей.};
\node[process,below=of f] (g) {Составьте уравнение. Решите его школьными методами с учётом ограничений.};
\node[process,below=of g] (h) {Проверьте скобки, знаки и размерности. Подставьте найденное значение в условие.};
\node[startstop,below=of h] (i) {Ответ с единицами};
\foreach \u/\v in {a/b,b/c,c/d,d/e,e/f,f/g,g/h,h/i}{\draw[arrow] (\u)--(\v);}
\end{tikzpicture}
\end{center}
\clearpage

\topic{7. Типичные ошибки и способы самопроверки}
\begin{enumerate}
\item \textbf{Неправильное выражение пути.} Если человек дошёл до опушки и вернулся на $x$ километров, его полный путь равен $L+x$. Если второй ещё не дошёл до опушки на $x$ километров, его путь равен $L-x$.
\item \textbf{Перепутаны большее и меньшее времена.} В задачи с общей точкой прибытия и разным временем старта раньше выехавший участник двигался дольше. Поэтому его время ставят первым при вычитании.
\item \textbf{Минуты оставлены вместо часов.} При скорости в километрах в час нужно записывать $8/60$ часа, а не $8$ часов.
\item \textbf{Потеря скобок.} Выражение $10x-10(x-5)$ равно $50$: знак минус меняет знаки у \emph{всех} слагаемых внутри скобок.
\item \textbf{Принят отрицательный корень для скорости.} Если скорость обозначена $x$, то $x>0$; если более медленная скорость $x-5$, то $x>5$.
\item \textbf{Ответ указан в других единицах.} Если задача спрашивает метры, после нахождения километров умножьте ответ на $1000$; если спрашивает минуты, после нахождения часов умножьте на $60$.
\end{enumerate}

\subtopic{Краткая памятка}
\begin{center}
\renewcommand{\arraystretch}{1.35}
\begin{tabularx}{\linewidth}{@{}p{0.43\linewidth}X@{}}
\toprule
Сигнал в условии & Что записать в уравнении \\
\midrule
Вышли одновременно из одного места; один дальше другого & Разность путей равна расстоянию между ними. \\
Один дошёл до конечной точки и повернул обратно & Полные пути: $L+x$ и $L-x$; времена одинаковы. \\
Один отправился раньше, но прибыли одновременно & Время раннего отправления больше на указанный промежуток. \\
Одно и то же расстояние, разные скорости & Сравнить $S/v_1$ и $S/v_2$. \\
Разность скоростей дана в метрах в минуту & Перевести её в километры в час до составления таблицы. \\
\bottomrule
\end{tabularx}
\end{center}

\clearpage
\topic{Тренировочный блок --- Путь А}
\textbf{10 задач.} Решайте последовательно, начиная с построения схемы там, где она нужна. Для задач с несколькими участниками рекомендуется составлять таблицу «скорость --- время --- расстояние». Все вычисления выполняйте самостоятельно; ответы помещены после блока упражнений.

\subtopic{Базовый уровень: одновременное движение}
\begin{enumerate}[label=\textbf{\arabic*.},wide=0pt]
\item Два пешехода одновременно вышли из одной и той же точки пешеходной аллеи и направились к городскому парку. Оба шли с постоянными скоростями, причём скорость первого пешехода была на $1{,}5\kmh$ больше скорости второго. Через какое количество минут после начала прогулки расстояние между пешеходами станет равным $450\met$? В ответе укажите время в минутах.

\item Два велосипедиста одновременно начали поездку от здания спортивной школы и поехали по одной прямой дороге в одном направлении. Первый велосипедист двигался с постоянной скоростью $16\kmh$, второй --- с постоянной скоростью $12\kmh$. Определите, какое расстояние в километрах будет между велосипедистами через $45$ минут после начала движения. Считайте, что оба велосипедиста всё это время ехали без остановок.

\item Два велосипедиста выехали одновременно с одной площадки и двигались в одном направлении по ровному участку дороги. Скорость первого велосипедиста составляла $17\kmh$, а скорость второго --- $13\kmh$. Через какое время после выезда расстояние между велосипедистами достигнет $6\kms$? Ответ выразите в минутах.
\end{enumerate}

\subtopic{Средний уровень: движение вдогонку с опозданием}
\begin{enumerate}[label=\textbf{\arabic*.},start=4,wide=0pt]
\item Бегун начал тренировку и выбежал со стадиона по прямой дорожке со скоростью $6\kmh$. Через $20$ минут после его старта с того же стадиона в том же направлении выехал велосипедист, скорость которого была постоянной и равнялась $18\kmh$. Через сколько минут после своего выезда велосипедист догонит бегуна? Предположите, что оба участника движения сохраняли указанные скорости и не останавливались.

\item Из небольшого посёлка по одной дороге выехал грузовой автомобиль и направился к складу со скоростью $35\kmh$. Спустя $24$ минуты с того же места в том же направлении отправился легковой автомобиль со скоростью $55\kmh$. Оба автомобиля ехали равномерно и без остановок. Через сколько минут после своего отправления легковой автомобиль догонит грузовой? Укажите время в минутах, отсчитываемое именно от выезда легкового автомобиля.
\end{enumerate}

\subtopic{Повышенный уровень: движение с разворотом}
\begin{enumerate}[label=\textbf{\arabic*.},start=6,wide=0pt]
\item Два туриста одновременно вышли от туристической базы в направлении смотровой площадки, расположенной в $6\kms$ от базы. Первый турист шёл со скоростью $9\kmh$, а второй --- со скоростью $6\kmh$. Первый турист дошёл до смотровой площадки, сразу развернулся и пошёл обратно по той же дороге. Через некоторое время он встретил второго туриста, который ещё продолжал идти к площадке. На каком расстоянии в метрах от смотровой площадки произошла встреча? Скорости туристов на всём протяжении пути считайте постоянными.

\item От входа в лесопарк два человека одновременно направились к беседке, находящейся в $8\kms$ от входа. Первый человек шёл со скоростью $12\kmh$, а второй --- со скоростью $8\kmh$. Более быстрый человек добрался до беседки, не задерживаясь повернул обратно и встретил второго человека на дороге между входом и беседкой. Найдите расстояние в метрах от беседки до места встречи. Считайте, что оба двигались равномерно и не делали остановок.
\end{enumerate}

\clearpage
\topic{Тренировочный блок --- продолжение}
\subtopic{Сложный уровень: разные моменты отправления}
\begin{enumerate}[label=\textbf{\arabic*.},start=8,wide=0pt]
\item Из города $A$ в город $B$, расстояние между которыми составляет $24\kms$, сначала выехал грузовой автомобиль. Через $6$ минут после его отправления по тому же маршруту выехал автобус. Автобус двигался с постоянной скоростью, на $8\kmh$ превышающей скорость грузового автомобиля. Несмотря на более поздний выезд, оба транспортных средства прибыли в город $B$ одновременно. Найдите скорость автобуса в километрах в час. Составьте уравнение, сравнив времена движения автомобилей.

\item Из деревни в соседний населённый пункт, находящийся на расстоянии $30\kms$, выехал первый мотоциклист. Через $18$ минут вслед за ним по тому же маршруту выехал второй мотоциклист, скорость которого была на $5\kmh$ больше скорости первого. Оба мотоциклиста прибыли в соседний населённый пункт в один момент времени. Определите скорость второго мотоциклиста в километрах в час. Считайте, что каждый из них двигался с постоянной скоростью без остановок.
\end{enumerate}

\subtopic{Сложный уровень: перевод скорости и уравнение}
\begin{enumerate}[label=\textbf{\arabic*.},start=10,wide=0pt]
\item Два поезда --- пассажирский и скорый --- преодолели один и тот же маршрут длиной $270\kms$. Пассажирский поезд двигался с постоянной скоростью, а скорый поезд за каждую минуту проходил на $300\met$ больше, чем пассажирский. В результате пассажирскому поезду на весь маршрут потребовалось на $1$ час $15$ минут больше, чем скорому. Найдите скорость пассажирского поезда в километрах в час. Перед составлением уравнения обязательно переведите разность скоростей в километры в час.
\end{enumerate}

\clearpage
\topic{Ответы к тренировочному блоку}
\textit{Сначала выполните задания самостоятельно, затем сверьте полученные результаты с таблицей.}
\vspace{5mm}
\begin{center}
\renewcommand{\arraystretch}{1.5}
\begin{tabularx}{0.86\linewidth}{@{}p{15mm}X@{}}
\toprule
№ & Верный ответ \\
\midrule
1 & $18$ минут \\
2 & $3$ километра \\
3 & $90$ минут \\
4 & $10$ минут после выезда велосипедиста \\
5 & $42$ минуты после выезда легкового автомобиля \\
6 & $1200$ метров от смотровой площадки \\
7 & $1600$ метров от беседки \\
8 & $48\kmh$ --- скорость автобуса \\
9 & $25\kmh$ --- скорость второго мотоциклиста \\
10 & $54\kmh$ --- скорость пассажирского поезда \\
\bottomrule
\end{tabularx}
\end{center}
\vspace{8mm}
\subtopic{Как проверить результат без полного повторения решения}
\begin{itemize}
\item \textbf{Задачи 1--3:} умножьте разность скоростей на найденное время в часах. Вы должны получить указанное расстояние между участниками.
\item \textbf{Задачи 4--5:} вычислите преимущество первого участника к моменту выезда второго. Разделите его на разность скоростей.
\item \textbf{Задачи 6--7:} найдите полные пути $L+x$ и $L-x$, а затем сравните времена их прохождения. Они должны совпасть.
\item \textbf{Задачи 8--9:} отдельно вычислите время движения каждого транспортного средства. Разница должна совпадать с указанным промежутком между отправлениями.
\item \textbf{Задача 10:} проверьте сначала разность скоростей, а затем разность времени прохождения маршрута.
\end{itemize}
\vfill
\begin{center}\small\color{mute}Лёвин Артём Александрович эксклюзивно для Ярослава Верещагина\end{center}
\end{document}
